\SourceLeaf{253}%
bestimmen sie die Koeffizienten \(\alpha\mathord{.}\ \beta,\ \gamma\mathord{.}\ \delta\). Es gilt
\[
\begin{aligned}
\alpha&=\BareRoot{\frac{E}{E'}}\mathord{.}\frac{\sin(\omega+\omega'-\psi)}{\sin\omega'},&\quad\beta&=\BareRoot{\frac{G}{E'}}\mathord{.}\frac{\sin(\omega'-\psi)}{\sin\omega'}\\
\gamma&=\BareRoot{\frac{E}{G'}}\mathord{.}\frac{\sin(\psi-\omega)}{\sin\omega'},&\delta&=\BareRoot{\frac{G}{G'}}\mathord{.}\frac{\sin\psi}{\sin\omega'}
\end{aligned}
\]
Hinzuzufügen sind die Gleichungen
\[
\begin{aligned}
\cos\omega&=\frac{F}{\BareRoot{EG}},&\cos\omega'&=\frac{F'}{\BareRoot{E'G'}},\\
\sin\omega&=\BareRoot{\frac{EG-FF}{EG}},&\sin\omega'&=\BareRoot{\frac{E'G'-F'F'}{E'G'}}
\end{aligned}
\]
sodass die vier Gleichungen auch so geschrieben werden können
\[
\begin{aligned}
\alpha\BareRoot{(E'G'-F'F')}&=\BareRoot{EG'}\mathord{.}\sin(\omega+\omega'-\psi)\\
\beta\BareRoot{(E'G'-F'F')}&=\BareRoot{GG'}\mathord{.}\sin(\omega'-\psi)\\
\gamma\BareRoot{(E'G'-F'F')}&=\BareRoot{EE'}\mathord{.}\sin(\psi-\omega)\\
\delta\BareRoot{(E'G'-F'F')}&=\BareRoot{GE'}\mathord{.}\sin\psi
\end{aligned}
\]

Da die Substitutionen \(\mathrm{d}p'=\alpha\mathrm{d}p+\beta\mathrm{d}q\) und \(\mathrm{d}q'=\gamma\mathrm{d}p+\delta\mathrm{d}q\) das Trinom \(E'\mathrm{d}p'^2+2F'\mathrm{d}p'\mathord{.}\mathrm{d}q'+G'\mathrm{d}q'^2\) in \(E\mathrm{d}p^2+2F\mathrm{d}p\mathord{.}\mathrm{d}q+G\mathrm{d}q^2\) überführen müssen, erhalten wir leicht
\[
EG-FF=(E'G'-F'F')(\alpha\delta-\beta\gamma)^2
\]
und da umgekehrt das letztere Trinom durch die Substitution
\[
(\alpha\delta-\beta\gamma)\mathrm{d}p=\delta\mathrm{d}p'-\beta\mathrm{d}q',\qquad(\alpha\delta-\beta\gamma)\mathrm{d}q=-\gamma\mathrm{d}p'+\alpha\mathrm{d}q'
\]
wieder in das erstere übergehen muss, finden wir
\[
\begin{aligned}
E\delta\delta-2F\gamma\delta+G\gamma\gamma&=\frac{EG-FF}{E'G'-F'F'}\mathord{.}E'\\
E\beta\delta-F(\alpha\delta+\beta\gamma)+G\alpha\gamma&=-\frac{EG-FF}{E'G'-F'F'}\mathord{.}F'\\
E\beta\beta-2F\alpha\beta+G\alpha\alpha&=\frac{EG-FF}{E'G'-F'F'}\mathord{.}G'
\end{aligned}
\]
\Article{22}
Von der allgemeinen Untersuchung des vorhergehenden Artikels gehen wir zu einer sehr weitreichenden Anwendung über. Dabei behalten \(p\) und \(q\) ihre allgemeinste Bedeutung, und für \(p',q'\) wählen wir die Größen, die in Artikel 15 mit \(r,\varphi\) bezeichnet wurden. Diese Zeichen
